\title{DeepSeekMath: Pushing the Limits of Mathematical Reasoning in Open Language Models}
\section{Supervised Fine-Tuning}

\begin{itemize}[topsep=0pt]
    \item \textbf{English mathematical datasets}: We enrich GSM8K and MATH with tool-augmented solution annotations, incorporate a portion of MathInstruct \citep{MathInstruct}, and use the training data from Lila-OOD \citep{lila}, all featuring solutions expressed using CoT or PoT frameworks. The English dataset covers a spectrum of mathematical disciplines, including but not limited to algebra, probability, number theory, calculus, and geometry.
    \item \textbf{Chinese mathematical datasets}: Our Chinese corpus features K-12 mathematics questions spanning 76 distinct subcategories such as linear equations. Solutions are presented in both CoT and tool-integrated reasoning styles.
\end{itemize}

Given these limitations, further research will be necessary and is left for subsequent investigations.

Nevertheless, these results should be interpreted with caution due to several caveats. Our investigation does not cover:

\item \textbf{Open-source models} comprise both general-purpose models and those with mathematics-specific enhancements. The general models are: (1) DeepSeek-LLM-Chat 67B \citep{deepseek-llm}, (2) Qwen 72B \citep{qwen}, (3) SeaLLM-v2 7B \citep{seallm}, and (4) ChatGLM3 6B \citep{chatglm3}. For models improved for mathematical reasoning, the collection includes: (5) InternLM2-Math 20B~\footnote{\url{https://github.com/InternLM/InternLM-Math}}, an extension of InternLM2 further trained on mathematical data and fine-tuned with instructions; (6) Math-Shepherd-Mistral 7B, which involves PPO training \citep{schulman2017proximal} applied to Mistral 7B \citep{mistral} using a process-level supervised reward model; (7) the WizardMath series \citep{wizardmath}, which enhances the mathematical reasoning capabilities of Mistral 7B and Llama-2 70B \citep{llama2} via evolve-instruct (a method of instruction tuning using instructions generated by AI) and PPO, primarily drawing from GSM8K and MATH datasets; (8) MetaMath 70B \citep{metamath}, which is a version of Llama-2 70B fine-tuned on an enriched set comprising GSM8K and MATH; (9) ToRA 34B \cite{tora}, which takes CodeLlama 34B and fine-tunes it for mathematical reasoning with integrated tools; and (10) MAmmoTH 70B \citep{MathInstruct}, which is an instruction-tuned variant of Llama-2 70B using MathInstruct.

Table \ref{tab:sft_rl_math} presents a comparative analysis of open-source and proprietary models using both chain-of-thought and tool-assisted reasoning approaches on English and Chinese evaluation tasks. The findings indicate: 1) When employing chain-of-thought reasoning, DeepSeekMath-RL 7B achieves 88.2\% accuracy on GSM8K and 51.7\% on MATH. These results exceed those attained by all open-source models within the 7B to 70B parameter range, and outperform most closed-source systems as well. 2) Notably, the training for DeepSeekMath-RL 7B is exclusively based on chain-of-thought instruction tuning data from GSM8K and MATH, using DeepSeekMath-Instruct 7B as its foundation. Despite its limited training data variety, DeepSeekMath-RL 7B consistently achieves higher scores than DeepSeekMath-Instruct 7B across all metrics, highlighting the substantial benefits of reinforcement learning.

\begin{itemize}[topsep=0pt]
    \item \textbf{Math Training for 150B Tokens}: DeepSeek-LLM 1.3B undergoes training directly on 150B tokens drawn from mathematics datasets;
    \item \textbf{Training on a mixture of 400B Code Tokens and 150B Math Tokens}: As sequential code and math training often diminishes coding ability, we also assess a scenario where the model is exposed to a joint dataset—interleaving 400B code tokens with 150B math tokens—to determine if this unified regimen both preserves code skills and enhances mathematical reasoning while addressing catastrophic forgetting.
\end{itemize}

\subsection{Lessons Learnt in Pre-Training}

We develop a comprehensive instruction-tuning dataset for mathematics that encompasses both English and Chinese language problems, drawing from a wide array of mathematical domains and levels of difficulty. Each problem is matched with a solution crafted in various reasoning formats: chain-of-thought (CoT) \citep{cot}, program-of-thought (PoT) \citep{pot,pal}, or with tool-based reasoning as proposed in \citep{tora}. Our compiled training set consists of 776,000 examples.

\subsubsection{How to Achieve More Effective RL?}
\label{sec:effec_rl}
Our results affirm that RL performs admirably on mathematical reasoning challenges. In addition, we outline a unifying perspective that frames various prominent training approaches as forms of direct or streamlined RL methods. According to Equation \ref{eq:objective}, three core elements structure this paradigm: Data Source, Algorithm, and Reward Function. We suggest possible future research trajectories for each element.

\subsection{Summary of Evaluations and Metrics}

\section{Conclusion, Limitation, and Future Work} 

Automating formal proofs is vital for improving the accuracy, trustworthiness, and efficiency of mathematical proof processes, which has garnered increasing research interest. We assess DeepSeekMath-Base 7B’s capabilities in informal-to-formal proof generation, following the setup in \citep{dsp_proof}: given an informal conjecture, its formal representation, and an informal proof, the model must output a formalized proof. Experiments are conducted on the miniF2F benchmark \citep{minif2f}, focused on Olympiad-level formal mathematics, requiring models to synthesize formal Isabelle proofs via few-shot prompting. In accordance with \cite{dsp_proof}, we instruct the models to produce preliminary proof outlines (sketches), after which the Sledgehammer automated prover \citep{sledgehammer} completes any unfinished proof details. Results in Table \ref{tab:base_math_pot_proof} reveal that DeepSeekMath-Base 7B delivers compelling performance in proof autoformalization.

    \item Through our experimentation in math pre-training, we observe that preceding math-focused training with code-related training enhances a model’s proficiency in solving mathematical problems, regardless of whether external tools are utilized. This finding provides empirical insight into the ongoing debate: \textit{does code training improve reasoning abilities?} For mathematical reasoning, our evidence supports an affirmative response.

\subsection{Contributions}
This work makes significant advancements in scalable mathematics-oriented pre-training, as well as an in-depth investigation of reinforcement learning methods.

To assess the quality of the DeepSeekMath~Corpus, we conduct pre-training experiments to compare it against other recent large-scale math-focused corpora:

\section{Math Pre-Training}

\paragraph{Results}
Performance metrics for each experimental configuration are reported in Table \ref{tab:code-math} and Table \ref{tab:code-math-general-eval}.

\paragraph{Reward Function} The reward function provides the foundational training feedback. In reinforcement learning, this is commonly instantiated as a neural reward model. We identify three critical research avenues for advancing reward models: 1) \textbf{Improving the generalization capacity of reward models.} To achieve genuine improvements in LLM performance, the reward model must successfully generalize to previously unseen questions and sophisticated decoding outputs; otherwise, RL might only maintain current distributions rather than advance underlying abilities; 2) \textbf{Quantifying reward model uncertainty.} Accurately modeling uncertainty may facilitate the integration between imperfect reward models and weak-to-strong algorithmic frameworks; 3) \textbf{Developing process-level reward models with efficiency and precision.} Such models are needed to supply granular feedback throughout the reasoning trajectory \citep{lightman2023let,wang2023math}.

Initially, the OpenWebMath dataset \citep{openwebmath}, known for its high-quality mathematical web documents, is selected as the seed corpus. We utilize this corpus to train a fastText classifier \citep{joulin2016fasttext} to identify and retrieve web pages similar to those in OpenWebMath. For this, we randomly sample 500,000 entries from the seed set as positive examples, while 500,000 randomly chosen web pages from Common Crawl serve as negative instances. The model is trained using an open-source fastText implementation\footnote{\url{https://fasttext.cc}}, with hyperparameters set as follows: 256-dimensional vectors, a learning rate of 0.1, a maximum word n-gram length of 3, a minimum word frequency of 3, and 3 training epochs. To compress the massive Common Crawl repository, we perform both URL-level deduplication and near-duplicate detection, narrowing down to 40B unique HTML pages. Utilizing the fastText classifier, we extract mathematical content from this refined set. Collected web pages are then ranked by their fastText prediction scores, and only the highest-scoring documents are retained as high-quality math data. The quantity of preserved data is further refined through pre-training experiments using subsets of the top 40B, 80B, 120B, and 160B tokens, with the initial iteration retaining the top 40B tokens.

Allowing models to utilize both natural language reasoning and code-based tools during problem solving, DeepSeekMath-Instruct 7B attains close to 60\% accuracy on the MATH benchmark, exceeding the performance of all previously released open-source models. In other benchmark assessments, our model demonstrates results on par with DeepSeek-LLM-Chat 67B, which is tenfold larger and was previously the best among open-source options.

The formulation features three primary elements: 1) \textit{Data Source $\mathcal{D}$}, which specifies the dataset used during training; 2) \textit{Reward Function $\pi_{{rf}}$}, which serves as the mechanism for generating reward signals to guide training; and 3) \textit{Algorithm $\mathcal{A}$}, responsible for integrating both training data and reward information to compute the gradient coefficient $GC$, determining the strength of the penalty or positive feedback on data samples. Leveraging this unified framework, we consider several widely adopted approaches:

Despite its strong results on quantitative reasoning datasets, DeepSeekMath~still lags behind closed models in tasks involving geometry and proof construction. As observed during preliminary evaluations, the model struggles with geometric problems such as those involving triangles or ellipses, possibly reflecting biases introduced during the data collection and curation phases. Additionally, due to its model size limitations, DeepSeekMath~does not match GPT-4 in few-shot learning: whereas GPT-4 benefits from few-shot examples, DeepSeekMath~exhibits nearly the same results whether evaluated in zero-shot or few-shot conditions. Moving forward, our plans include refining the data selection methodology to compile higher-quality pre-training datasets. We will also further investigate the directions outlined in Section \ref{sec:effec_rl} to achieve more effective reinforcement learning strategies for large language models.

\subsection{Data Collection and Decontamination}
This section describes the methodology used for assembling the DeepSeekMath Corpus from Common Crawl data. As illustrated in Figure \ref{fig:data_collect}, we developed an iterative framework for systematically extracting a large-scale mathematical dataset, beginning with a foundational “seed” corpus consisting of a select set of high-quality math-focused resources. Notably, the outlined procedure can be readily extended to other specialized fields such as programming.

\subsection{Training and Evaluating DeepSeekMath-Base 7B}

\subsubsection{From PPO to GRPO} The actor-critic approach Proximal Policy Optimization (PPO) \citep{schulman2017proximal} is a standard method employed in the reinforcement learning fine-tuning of LLMs \citep{ouyang2022training}. This algorithm aims to improve LLMs by optimizing the following surrogate loss:

Here, $\varepsilon$ and $\beta$ denote hyper-parameters, and $\hat{A}_{i,t}$ refers to the advantage computed exclusively using the relative rewards among the outputs in each group—details on this computation follow in subsequent subsections. The methodology underlying GRPO’s advantage computation is well aligned with the inherently comparative training of reward models, which often rely on evaluating differences between multiple responses to identical questions. It is also noteworthy that, distinct from standard practice where the KL penalty is incorporated within the reward, GRPO regularizes the optimization objective directly by penalizing the KL divergence between the new policy and the reference policy—thereby keeping the computation of $\hat{A}_{i,t}$ more straightforward. Furthermore, as opposed to the KL penalty definition in (\ref{eq:PPO-reward}), we

\begin{equation}
    \nabla_{\theta}\mathcal{J}_{\textcolor{red}{\mathcal{A}}}(\theta) = \mathbb{E}[\underbrace{(q,o) \sim \textcolor{red}{\mathcal{D}}}_{Data \ Source}]\left( \frac{1}{|o|} \sum_{t=1}^{|o|}  \underbrace{GC_{{\mathcal{A}}}(q, o, t, \textcolor{red}{\pi_{{rf}}})}_{Gradient \ Coefficient}  \nabla_{\theta}\log \pi_{\theta}(o_t | q, o_{<t})\right).
\label{eq:objective}
\end{equation}

To gauge mathematical reasoning capabilities, we assess model performance with and without tool assistance, leveraging four quantitative reasoning evaluation suites in both English and Chinese. Our methodical evaluation involves direct comparison with other leading models of the era:

As detailed in Table \ref{tab:base_math_cot}, DeepSeekMath-Base 7B consistently achieves top scores across all eight evaluation benchmarks for open-source base models, outperforming prominent options such as Mistral 7B \citep{mistral} and the newly introduced Llemma 34B \citep{llemma}, the latter of which benefits from specialized training on Proof-Pile-2 \citep{llemma}. Particularly on the challenging MATH benchmark, which targets competition-level problems, DeepSeekMath-Base achieves an absolute margin of over 10\% above the best-performing open-source base models to date, and even surpasses Minerva 540B \citep{minerva}—a closed-source model built on PaLM \citep{palm}, 77 times larger in scale and specifically fine-tuned on mathematical materials.

\begin{itemize}[topsep=0pt]
    \item \textbf{MathPile} \citep{mathpile}: a diverse collection of sources totaling 8.9B tokens, aggregated from textbooks, Wikipedia, ProofWiki, CommonCrawl, StackExchange, and arXiv—with arXiv accounting for more than 85\% of the dataset;
    \item \textbf{OpenWebMath} \citep{openwebmath}: derived from CommonCrawl, containing 13.6B tokens of math-filtered content;
    \item \textbf{Proof-Pile-2} \citep{llemma}: a compilation comprising OpenWebMath, AlgebraicStack (10.3B tokens of code), and arXiv papers (28.0B tokens). In our Proof-Pile-2 experiments, we maintain the arXiv:Web:Code token ratio of 2:4:1 as recommended in \cite{llemma}.
\end{itemize}

For evaluation, DeepSeek-LLM 1.3B was trained on 150B tokens and DeepSeek-Coder-Base-v1.5 7B on 40B tokens from each arXiv-derived corpus individually. The results consistently indicate that exclusive training on arXiv datasets fails to yield significant enhancement—and occasionally results in poorer outcomes—across a spectrum of mathematical benchmarks with varying levels of complexity. These benchmarks encompass quantitative reasoning sets like GSM8K and MATH (see Table \ref{tab:arxiv-cot}), multiple-choice tasks such as MMLU-STEM (Table \ref{tab:arxiv-cot}), and formal mathematics benchmarks including miniF2F (Table \ref{tab:arxiv_atp}).

\subsubsection{Training Setting}
\label{sec:quality-policy}
We conduct mathematical training using a general-purpose language model with 1.3B parameters, which adheres to the architecture of the DeepSeek LLMs \citep{deepseek-llm}, henceforth referred to as DeepSeek-LLM 1.3B. Individual models are trained from scratch on each mathematical dataset for a total of 150B tokens. All training procedures utilize the lightweight and efficient HAI-LLM framework \citep{haillm}. In alignment with the DeepSeek LLMs training methodology, the AdamW optimizer \citep{adamW} is used, setting $\beta_1=0.9$, $\beta_2=0.95$, and $\mathrm{weight\_decay}=0.1$. The learning rate follows a multi-step scheduling: it increases linearly, reaching its maximum value after 2,000 steps of warmup, then decays to 31.6\% of its peak at 80\% of total training, and further declines to 10.0\% of the peak by the 90\% completion mark. The highest learning rate is set at 5.3e-4, with training performed using batches of 4M tokens and a maximum context window of 4K tokens.

Figure \ref{fig:rl-analysis} illustrates that Online RFT demonstrates a marked improvement over RFT on both evaluation benchmarks. In particular, during early training stages, Online RFT and RFT yield similar performance; however, as training progresses, Online RFT attains a clear and consistent advantage, evidencing the benefits of employing online training protocols. This aligns with expectations: initially, the actor and SFT models behave similarly, and discrepancies in their sampled outputs are minimal. Over time, as the actor diverges from the SFT model, the distinctions in its sampled data grow, enhancing the advantages conferred by real-time data collection.

In this work, we introduce DeepSeekMath, which surpasses all existing open-source models on the competition-level MATH benchmark and nearly matches the results of proprietary models. DeepSeekMath~builds on DeepSeek-Coder-v1.5 7B as its foundation and is continually trained over 500B tokens, including a substantial subset of 120B mathematical tokens extracted from Common Crawl. Our comprehensive ablation experiments reveal that online resources provide considerable value as sources of quality mathematical data, whereas arXiv contributes less than anticipated. We propose Group Relative Policy Optimization (GRPO), a new adaptation of Proximal Policy Optimization (PPO), which can substantially enhance mathematical reasoning with reduced memory requirements. Experimental analysis indicates that GRPO remains effective, even when DeepSeekMath-Instruct 7B attains high performance on established benchmarks. Furthermore, we present a unifying framework to contextualize multiple methods, and we identify promising avenues for advancing reinforcement learning.

    \item The DeepSeekMath-Base 7B, our foundational pre-trained model, delivers performance on par with Minerva 540B \citep{minerva}, suggesting that model size alone does not solely determine mathematical reasoning prowess. These results underscore that smaller architectures, when trained on superior quality data, can also achieve robust mathematical competence.

These models are generally designed for broad application and typically have been subjected to various alignment techniques.

\item \textbf{Natural Language Understanding, Reasoning, and Code}: To thoroughly evaluate the model’s general abilities in language comprehension, reasoning, and programming, we benchmark DeepSeekMath-Base with Massive Multitask Language Understanding (MMLU) \citep{mmlu}, which includes 57 multiple-choice tasks from diverse domains, and BIG-Bench Hard (BBH) \citep{bbh}, comprising 23 complex problems requiring multi-step reasoning. Additionally, HumanEval \citep{codex} and MBPP \citep{mbpp}, commonly used in assessing code generation models, are used for further evaluation. Pre-training on mathematical data yields improvements in both language understanding and reasoning capabilities.

\paragraph{Data Source} The data source forms the foundation of every training procedure. For RL, we define it as the pool of unlabeled queries coupled with candidate responses sampled from the policy model. This work limits the data source to the questions used for instruction tuning and employs a basic nucleus sampling approach. We posit that this limitation may contribute to our RL pipeline’s primary impact on Maj@K improvement. Future investigations will assess RL pipelines with out-of-distribution prompts, together with \textbf{advanced sampling (decoding) strategies}, including those leveraging tree-search methodologies \citep{yao2023tree}. Furthermore, \textbf{efficient inference techniques} \citep{xia-etal-2023-speculative,leviathan2023fast,kwon2023efficient,xia2024unlocking}, which are crucial for the effective exploration by policy models, remain of paramount significance.

\section{Reinforcement Learning}

\begin{itemize}
    \item  \textbf{Supervised Fine-tuning (SFT)}: This method adapts a pretrained model using curated, human-annotated SFT datasets through further supervised training.
    \item \textbf{Rejection Sampling Fine-tuning (RFT)}: RFT conducts additional fine-tuning of the SFT model by employing outputs sampled from the SFT model, conditioned on the same prompts, and retaining only those outputs which are deemed correct.
    \item \textbf{Direct Preference Optimization (DPO)}: DPO enhances the SFT model by performing fine-tuning with extra model-generated responses, applying a pairwise DPO objective.
    \item \textbf{Online Rejection Sampling Fine-tuning (Online RFT)}: Unlike traditional RFT, this variant begins with the SFT model but updates the policy through fine-tuning on augmented outputs sampled directly from the evolving policy in real-time.
    \item \textbf{PPO/GRPO}: These methods initialize the policy from the SFT model, then optimize it by sampling outputs from the live policy model and providing reinforcement.
\end{itemize}

\subsubsection{ArXiv Papers Appear Limited in Enhancing Mathematical Reasoning}

\subsection{Group Relative Policy Optimization} Previous work has established that reinforcement learning (RL) methods are instrumental in further advancing mathematical reasoning capabilities in large language models (LLMs) after initial Supervised Fine-Tuning (SFT) \citep{wang2023math,wizardmath}. Here, we present a new RL technique called Group Relative Policy Optimization (GRPO), designed to be both efficient and effective.

\paragraph{Formal Mathematics}

    On these English benchmarks, DeepSeekMath-Base attains performance on par with the proprietary Minerva 540B \citep{minerva}, outperforming all existing open-source base models—including Mistral 7B \citep{mistral} and Llemma-34B \citep{llemma}—often by wide margins, independent of prior mathematical pre-training. For Chinese mathematical tasks, DeepSeekMath-Base consistently leads, which can be attributed to the inclusion of curated non-English datasets, as opposed to previous works \citep{minerva,llemma} which used predominantly English math corpora. After mathematical instruction tuning and the application of reinforcement learning, DeepSeekMath-Instruct and DeepSeekMath-RL display strong capabilities, achieving over 50\% accuracy on the challenging MATH benchmark—the first such result achieved by an open-source model.

\begin{itemize}[topsep=0pt]
    \item The potential value of arXiv-derived tokens for targeted mathematical use cases outside the scope of our work, such as converting formal proofs into informal language;
    \item The influence of incorporating arXiv tokens alongside other training data sources;
    \item Whether scaling up the model size would reveal benefits from training on arXiv content.
\end{itemize}

\subsubsection{Code Training Benefits Mathematical Reasoning}

\subsubsection{Process Supervision RL with GRPO} While outcome supervision provides rewards only after entire responses are produced, this feedback may be inadequate for effectively guiding learning, especially for intricate mathematical problems. Inspired by \cite{wang2023math}, we also implement process supervision, whereby a reward is assigned following each step in the reasoning trajectory. Specifically, for a prompt $q$ and $G$ generated responses $\{o_1, o_2, \cdots, o_G\}$, a process reward model allocates a reward to every reasoning step, resulting in: $\mathbf{R} = \{ \{r_1^{index(1)},\cdots,r_1^{index(K_1)}\}, \cdots,  \{r_G^{index(1)},\cdots,r_G^{index(K_G)}\} \}$, where $index(j)$ denotes the end position of the $j$-th reasoning step, and $K_i$ is the number of steps in the $i$-th response. These step-level rewards are group-normalized as $\widetilde{r}_i^{index(j)} = \frac{r_i^{index(j)} - {\rm mean(\mathbf{R})}}{{\rm std(\mathbf{R})}}$.
After normalization, the advantage at token $t$ is the sum of normalized future step rewards: $\hat{A}_{i, t} = \sum_{index(j) \ge t} \widetilde{r}_i^{index(j)}$.
The policy is then updated to maximize the objective specified

\begin{itemize}[topsep=0pt]
    \item We demonstrate that Common Crawl, an openly available internet corpus, holds substantial resources for mathematical applications. Leveraging a carefully crafted data filtering system, we compile the DeepSeekMath Corpus, a rigorously curated dataset containing 120B tokens sourced from web content selected for mathematical relevance. This dataset surpasses the math-specific web data utilized by Minerva \citep{minerva} by nearly a factor of 7, and exceeds the scale of OpenWebMath \citep{openwebmath} by 9 times.

Additionally, we develop a unified framework for understanding a range of methodologies, such as Rejection Sampling Fine-Tuning (RFT) \citep{yuan2023scaling}, Direct Preference Optimization (DPO) \citep{dpo}, PPO, and GRPO. This unified perspective reveals that all these approaches can be viewed as forms of direct or streamlined reinforcement learning strategies. We conduct a comprehensive set of studies—including comparisons of online vs. offline learning, outcome-based vs. process-based supervision, and single-turn vs. iterative RL—to explore key factors influencing this paradigm. Finally, we provide insights into how reinforcement learning enhances the capabilities of instruction-tuned models and propose further research avenues for more effective RL within this integrated paradigm.

\begin{equation}
\footnotesize
            \mathcal{J}_{PPO}(\theta) = \mathbb{E}{[q \sim P(Q), o \sim \pi_{\theta_{old}}(O|q)]} \frac{1}{|o|} \sum_{t=1}^{|o|} \min \left[ \frac{\pi_\theta(o_{t} | q, o_{<t})}{\pi_{\theta_{old}}(o_{t} | q, o_{<t})} A_{t}, \text{clip} \left( \frac{\pi_\theta(o_{t} | q, o_{<t})}{\pi_{\theta_{old}}(o_{t} | q, o_{<t})}, 1 - \varepsilon, 1 + \varepsilon \right)  A_{t} \right] ,
\end{equation}

Additionally, code-centric pretraining enhances mathematical problem-solving even without recourse to external tools. Sequential training where code data precedes math data leads to incremental improvements from the outset and streamlines subsequent adaptation to math tasks, culminating in peak results. However, an integrated training corpus blending code and math reduces performance for math problems solved without tools. A plausible explanation is that the restricted parameter count in DeepSeek-LLM 1.3B hinders the model’s ability to simultaneously internalize the nuances of both code and mathematical domains when trained together.

\begin{itemize}[topsep=0pt]
    \item \textbf{Code Training for 400B Tokens $\rightarrow$ Math Training for 150B Tokens}: We first pre-train DeepSeek-LLM 1.3B on 400B code tokens, followed by an additional 150B tokens consisting of mathematics content;
    \item \textbf{General Training for 400B Tokens $\rightarrow$ Math Training for 150B Tokens}: For comparison, we repeat the training protocol using 400B tokens sampled from a large-scale general corpus provided by DeepSeek-AI as the first-stage data instead of code, aiming to analyze the specific contributions of code data versus general data in strengthening mathematical reasoning.
\end{itemize}

    \item Contrary to common practice—especially among works focused on mathematics—we find that incorporating arXiv papers into training regimes does not yield measurable improvements across any of the mathematical benchmarks considered in our study.
\end{itemize}

\paragraph{Mathematical Problem Solving with Step-by-Step Reasoning}

\paragraph{Mathematical Problem Solving with Tool Use} To examine the model’s proficiency in program-assisted mathematical reasoning, we apply few-shot program-of-thought prompting \citep{pot,pal} on GSM8K and MATH datasets. For each task, the models generate Python scripts that leverage tools like \textit{math} and \textit{sympy} for complex operations; the computed outcome is treated as the problem’s solution. As reported in Table \ref{tab:base_math_pot_proof}, DeepSeekMath-Base 7B establishes a new performance benchmark, exceeding the results achieved by the previously leading Llemma 34B.

\begin{itemize}[topsep=0pt]
    \item \textbf{MathPile} \citep{mathpile}: A curated 8.9B-token dataset employing a set of heuristic cleaning and filtering procedures, comprising over 85\% scientific arXiv articles;
    \item \textbf{ArXiv-RedPajama} \citep{redpajama}: An entire collection of arXiv LaTeX documents that have been stripped of preambles, comments, macros, and bibliography, totaling 28.0B tokens.
\end{itemize}

Following the initial round of data collection, a significant number of mathematical web pages remained unacquired. This shortfall primarily arises from the fastText model's reliance on a set of positive training samples that are not sufficiently diverse. To address this limitation, we sought to broaden the coverage of the seed corpus by incorporating additional mathematical web resources, thereby improving the fastText model's performance. Our method begins with partitioning the entire Common Crawl dataset into non-overlapping domains, where each domain is defined by web pages under the same base URL. We then determine, for each domain, what proportion of its pages were gathered during the first data collection round. Domains exhibiting a collection rate greater than 10\% are labeled as math-oriented (e.g., \url{mathoverflow.net}).

\begin{itemize}[topsep=0pt]
    \item We propose Group Relative Policy Optimization (GRPO), a novel and resource-efficient reinforcement learning method. Unlike Proximal Policy Optimization (PPO), GRPO does not utilize a critic model; instead, it derives the baseline using group scores, which substantially lowers the computational and training cost.
    \item Our experiments reveal that GRPO considerably boosts the effectiveness of our instruction-tuned model, DeepSeekMath-Instruct, relying exclusively on instruction-tuning datasets. In addition, we note consistent improvements on out-of-domain tasks during reinforcement learning.
    \item We introduce a unified analytical framework to relate approaches such as RFT, DPO, PPO, and GRPO. Additionally, we perform extensive empirical analyses—examining dimensions such as online versus offline training, outcome versus process supervision, and single-turn versus iterative reinforcement learning—to systematically probe the core components of this framework.
    \item Leveraging this unified perspective, we investigate the mechanisms driving reinforcement learning’s success, and outline several promising avenues to further improve reinforcement learning for large language models.
\end{itemize}

We also assess natural language comprehension on MMLU \citep{mmlu}, reasoning capabilities via BBH \citep{bbh}, and programming proficiency on HumanEval \citep{codex} and MBPP \citep{mbpp}. Table \ref{tab:understanding_reasoning_code} demonstrates that DeepSeekMath-Base 7B achieves substantial gains over its predecessor, DeepSeek-Coder-Base-v1.5 \citep{deepseek-coder}, in both MMLU and BBH, underscoring the advantageous effect of mathematical pretraining on downstream language understanding and logical inference. Moreover, by incorporating code-specific tokens in ongoing training, DeepSeekMath-Base 7B preserves the coding performance observed in DeepSeek-Coder-Base-v1.5 on both programming benchmarks. Collectively, DeepSeekMath-Base 7B exhibits superior results relative to the general-purpose Mistral 7B model \citep{mistral} across all evaluated tasks in reasoning and code generation.

This work presents DeepSeekMath, a specialized language model designed to excel at mathematical tasks, surpassing the abilities of other open-source models and coming close to the benchmark performance of GPT-4 in academic evaluations. Central to our approach is the construction of the DeepSeekMath Corpus, an extensive, high-quality pre-training dataset containing 120B mathematical tokens. We compile this corpus from Common Crawl (CC) utilizing a classifier based on fastText \citep{joulin2016fasttext}. Initially, the classifier is trained with positive samples drawn from OpenWebMath \citep{openwebmath} and a variety of other webpages as negative samples. The trained classifier is subsequently employed to extract more positive samples from CC, which undergo further quality control through human annotation. The classifier is iteratively retrained using this improved dataset, resulting in enhanced classification accuracy. Empirical evaluation demonstrates that the large-scale corpus is of high caliber, as evidenced by DeepSeekMath-Base 7B attaining 64.2\% on GSM8K \citep{gsm8k} and 36.2\% on the MATH competition dataset \citep{MATH}, both surpassing the performance of Minerva 540B \citep{minerva}. Notably, since the DeepSeekMath Corpus is multilingual, we observe performance gains on Chinese mathematical benchmarks \citep{wei2023cmath,agieval} as well. We view our methodology for processing mathematical data as an initial foundation for further inquiry, highlighting considerable potential for subsequent advancements.

\textbf{Exploration and Analysis of Reinforcement Learning}

\end{document}

While arXiv papers frequently comprise a substantial portion of mathematical pre-training datasets \citep{minerva,gpt-f,llemma,mathpile}, there remains little systematic analysis of their actual contribution to models’ mathematical reasoning abilities. Contrary to common expectation, our findings suggest that arXiv papers provide minimal gains for mathematical reasoning capabilities. We tested this hypothesis using models with differing parameter counts, specifically DeepSeek-LLM 1.3B and DeepSeek-Coder-Base-v1.5 7B \citep{deepseek-coder}, and considered arXiv datasets processed through different workflows:

In this context, $r_\varphi$ indicates the reward model, $\pi_{ref}$ specifies the reference model, typically chosen as the initial SFT model, and $\beta$ serves as the scaling parameter for the KL penalty.

\subsection{Training and Evaluating DeepSeekMath-Instruct 7B}

Here, $\pi_{\theta}$ and $\pi_{\theta_{old}}$ represent the present and previous versions of the policy models, while $q$ and $o$ denote questions and outputs that are sampled respectively from the question dataset and from the output distribution under $\pi_{\theta_{old}}$. The term $\varepsilon$ stands for a clipping hyper-parameter introduced in PPO, aiming to promote training stability. The advantage term, $A_t$, is determined by Generalized Advantage Estimation (GAE) \citep{gae}, which relies on future rewards $\{r_{\ge t}\}$ and a separately learned value function $V_{\psi}$. In PPO, this necessitates joint optimization of both the policy and value function models. To prevent the reward model from being excessively optimized, a common mitigation is the introduction of a per-token KL penalty against a reference model, computed at each token during reward calculation \citep{ouyang2022training}, expressed as:

\begin{equation}
    r_{t} = r_\varphi(q, o_{\le t}) - \beta \log\frac{\pi_{\theta}(o_{t}|q, o_{<t})}{\pi_{ref}(o_{t}|q, o_{<t})},
\label{eq:PPO-reward}
\end{equation}

We proceed by manually tagging the URLs within these math-related domains that host mathematical content (such as \url{mathoverflow.net/questions}). Those pages, associated with these URLs but not previously collected, are subsequently incorporated into the seed corpus. This process yields a richer set of positive samples for retraining the fastText model, enabling improved retrieval of mathematical data in future iterations. After completing four cycles of this procedure, we amassed a total of 35.5M web pages related to mathematics, comprising approximately 120B tokens. In the fourth cycle, it became apparent that close to 98\% of the content had already been acquired during the third iteration, leading us to terminate further data collection at that stage.

\subsubsection{Iterative RL with GRPO} During the progression of reinforcement learning training, previously constructed reward models may no longer be adequate for supervising the evolving policy model. To address this, we implement iterative RL with GRPO. As detailed in Algorithm \ref{alg:iter-grpo}, iterative GRPO involves producing new datasets for the reward model by sampling outputs from the latest policy model, and further refining the reward model with a replay buffer that includes 10\% of earlier training data. Subsequently, we designate the policy model as the new reference model and continue to update it, utilizing the improved reward model for guidance.

A widely held but inadequately tested conjecture is that exposure to code enhances a model’s reasoning skills. Here, we provide evidence supporting this hypothesis in the context of mathematics: training on code improves the model’s capabilities in mathematical reasoning, regardless of whether tools are involved.

\paragraph{Natural Language Understanding, Reasoning, and Code}

Since the value function used within PPO generally constitutes a neural network similar in size to the policy model, its inclusion significantly increases memory usage and computational expense. Moreover, during reinforcement learning, the value function operates as a baseline to reduce variance when computing advantages. In large language model settings, the reward model is commonly configured to supply a reward for only the final token in the generated sequence, which presents challenges for learning an accurate, token-level value function. To overcome these limitations, we introduce Group Relative Policy Optimization (GRPO), as illustrated in Figure \ref{fig:grpo}. GRPO eliminates the requirement for value function approximation—central to PPO—and instead employs the mean reward across several output samples (generated for the same question) as the baseline. Concretely, for each question $q$, a set of outputs $\{o_1, o_2, \cdots, o_G\}$ is drawn from the prior policy $\pi_{\theta_{old}}$. The policy is subsequently refined by maximizing the objective below:

\textbf{Two-Stage Training}

\begin{itemize}[topsep=0pt]
    \item \textbf{High-quality}: 
    To assess downstream capabilities, we measure performance across 8 mathematical benchmarks employing few-shot chain-of-thought prompting \cite{cot}. Results in Table \ref{tab:corpora_comparison} demonstrate a notable superiority of models trained on the DeepSeekMath~Corpus. Furthermore, Figure \ref{fig:corpora_comparisons} reveals that training on DeepSeekMath~Corpus results in better accuracy than Proof-Pile-2 after 50B tokens (equivalent to one epoch over Proof-Pile-2), suggesting a higher average data quality in DeepSeekMath.
    \item \textbf{Multilingual}:
    DeepSeekMath~Corpus incorporates multilingual mathematical texts, with English and Chinese forming the majority of entries. Table \ref{tab:corpora_comparison} shows that models trained on DeepSeekMath~Corpus display enhanced mathematical reasoning proficiency in both English and Chinese, whereas prior datasets, which are primarily focused on English, show minimal to negative gains for Chinese tasks.
    \item \textbf{Large-scale}: 
    DeepSeekMath~Corpus surpasses prior mathematical datasets by several multiples in terms of size. As seen in Figure \ref{fig:corpora_comparisons}, DeepSeek-LLM 1.3B models trained with DeepSeekMath~Corpus not only learn more rapidly, but also continue to make progress as training proceeds. Conversely, baseline corpora are limited in scale, have been subject to multiple passes during training, and show early saturation in downstream performance.
\end{itemize}

Across both staged and joint-training paradigms, pre-training with code tokens confers clear benefits to models tackling program-augmented mathematical tasks. Notably, in the two-stage setup, initial code-focused training substantially augments accuracy on benchmarks such as GSM8K and MATH when using Python. Introducing a subsequent math-specific phase further elevates model capability. Furthermore, in the one-stage configuration, simultaneous exposure to code and math tokens curtails the degradation in code proficiency observed with strict sequential training, and amplifies performance for both general coding (Table \ref{tab:code-math-general-eval}) and program-mediated mathematical reasoning (Table \ref{tab:code-math}).

\section{Discussion}
This section elaborates on insights gained from both pre-training and reinforcement learning experiments.

\subsection{Training and Evaluating DeepSeekMath-RL}

\section{Introduction}

\subsection{SFT Data Curation}

This section presents DeepSeekMath-Base 7B, a foundational model designed for robust reasoning, particularly within mathematical domains. The initialization of our model is based on DeepSeek-Coder-Base-v1.5 7B \citep{deepseek-coder}, followed by training over a corpus of 500B tokens. The training data consists of 56\% from the DeepSeekMath Corpus, 4\% from AlgebraicStack, 10\% sourced from arXiv, 20\% from Github code repositories, while the final 10\% comprises natural language text from Common Crawl, available in both English and Chinese. The training configuration closely mirrors that outlined in Section \ref{sec:quality-policy}, with modifications including setting the peak learning rate to 4.2e-4 and utilizing a token batch size of 10M.

\subsubsection{Why RL Works?} This study applies reinforcement learning to a selected subset of instruction tuning data, resulting in notable gains over the base instruction tuning model. To shed light on the mechanism behind this improvement, we compare the Pass@K and Maj@K metrics for both Instruct and RL models across two evaluation benchmarks. Figure \ref{fig:maj-pass} displays that RL brings marked improvement to Maj@K scores, while Pass@K remains largely unchanged. This suggests that RL strengthens the overall output robustness—specifically, \textbf{the observed gains arise from a greater likelihood of correct answers appearing within the TopK outputs, rather than enhancements to core model competency.} Likewise, \citep{wang2023making} report a \textbf{misalignment problem} in reasoning scenarios with the SFT model, and demonstrate that reasoning accuracy in SFT models can be elevated through preference alignment strategies \citep{yuan2023rrhf,song2023preference,wang2023making}.

As illustrated in Figure \ref{fig:rl-analysis}, GRPO outperforms online RFT, underscoring the effectiveness of modifying positive and negative gradient coefficients. Moreover, GRPO+PS achieves better outcomes than GRPO+OS, which demonstrates the advantage of employing step-aware, fine-grained gradient coefficients. Additionally, our analysis extends to iterative RL, where we perform two cycles of iteration in our experiments. Figure \ref{fig:iter-rl} reveals that this iterative process markedly enhances performance, particularly during the initial iteration.

To examine the influence of code training on mathematical reasoning, we evaluated both a two-stage training process and a single-stage training scheme:


\begin{abstract}
Mathematical reasoning remains a particularly demanding task for language models, given the high degree of structure and complexity involved. This study presents DeepSeekMath~7B, developed by further pre-training DeepSeek-Coder-Base-v1.5 7B on a corpus of 120B math-centric tokens, extracted from Common Crawl along with additional natural language and programming data. Remarkably, DeepSeekMath~7B attains a score of 51.7\% on the MATH benchmark—designed to assess competition-level mathematical ability—without dependence on external toolkits or ensemble voting methods, narrowing the gap to leading proprietary models like Gemini-Ultra and GPT-4. Using self-consistency with 64 samples, DeepSeekMath~7B's score rises to 60.9\% on the same benchmark. The superior performance of DeepSeekMath~is largely attributed to two central components: (1) leveraging a substantial volume of publicly accessible web data through a rigorously constructed data filtration and selection process, and (2) the implementation of Group Relative Policy Optimization (GRPO), an adaptation of Proximal Policy Optimization (PPO), which not only strengthens mathematical reasoning but also improves the memory efficiency during PPO-based training.
\end{abstract}

\textbf{One-Stage Training}

Our evaluation strategy thoroughly examines the mathematical reasoning abilities of DeepSeekMath-Base 7B, emphasizing its proficiency in independently generating comprehensive solutions, employing tools for problem-solving, and engaging in rigorous formal theorem proving. In addition to mathematical competencies, we also report the model’s general capabilities, such as natural language understanding, logical reasoning, and programming performance.

\paragraph{Algorithms} Algorithms leverage both the data and the reward signal, which together determine the gradient coefficients used to update model parameters. According to Equation \ref{eq:objective}, most contemporary approaches fundamentally \textbf{TRUST} the reward function, adjusting the conditional probability of specific tokens upward or downward in response. Nevertheless, it cannot be guaranteed that the reward signal remains trustworthy across all scenarios, particularly in highly intricate tasks. For instance, the PRM800K datasets \citep{lightman2023let}, despite being annotated by skilled professionals, still exhibit roughly 20\% annotation errors\footnote{\url{https://github.com/openai/prm800k/issues/12\#issuecomment-1728491852}}. This motivates an investigation into reinforcement learning algorithms that can tolerate and remain effective under noisy reward conditions. Our perspective is that \textbf{WEAK-TO-STRONG} \citep{burns2023weak} alignment techniques have the potential to fundamentally transform learning paradigms.

\subsubsection{Outcome Supervision RL with GRPO} In the case of outcome supervision, for a given prompt $q$, a set of responses $\{o_1, o_2, \cdots, o_G\}$ are sampled from the previous policy $\pi_{\theta_{old}}$. Each generated response is assigned a reward by a reward model, producing a vector of $G$ rewards $\mathbf{r}=\{r_1, r_2, \cdots, r_G\}$. These scores are then normalized within the group by subtracting their mean and dividing by their standard deviation. Under outcome supervision, this normalized score is used as a terminal reward for each output $o_i$, and the advantage $\hat{A}_{i, t}$ assigned to every token position $t$ in $o_i$ equals this normalized reward: $\hat{A}_{i, t} = \widetilde{r}_i = \frac{r_i- {\rm mean}(\mathbf{r})}{{\rm std}(\mathbf{r})}$. Policy updates are then performed by maximizing the target defined in equation (\ref{eq:GRPO-obj}).

\subsubsection{Evaluation Results}

\subsection{Validating the Quality of the DeepSeekMath~Corpus}

As detailed in Table \ref{tab:sft_rl_math}, when tool-based approaches are not allowed during evaluation, DeepSeekMath-Instruct 7B achieves high accuracy in multi-step reasoning. Notably, for the competitive MATH dataset, this model outperforms every open-source alternative and most commercial proprietary models (such as Inflection-2 and Gemini Pro), achieving a margin of at least 9 percentage points. This superiority holds even when compared to considerably larger models (e.g., Qwen 72B) or those augmented specifically via reinforcement learning focused on mathematical capabilities (for instance, WizardMath-v1.1 7B). Although DeepSeekMath-Instruct is competitive with proprietary Chinese systems like GLM-4 and Baichuan-3 on MATH, it has yet to match the performance of GPT-4 or Gemini Ultra.

\textbf{Math Pre-Training at Scale}

DeepSeekMath-Base is built upon DeepSeek-Coder-Base-v1.5 7B \citep{deepseek-coder}, since initializing from a code-centric model proves more effective than using a general-purpose LLM. Additionally, our findings indicate that math-oriented training not only advances the model's mathematical proficiency, but also enhances its general reasoning abilities, as demonstrated by improvements on the MMLU \citep{mmlu} and BBH \citep{bbh} evaluation benchmarks.

\begin{itemize}[topsep=0pt]
    \item \textbf{English and Chinese Mathematical Reasoning}: Our evaluation spans a wide array of English and Chinese mathematical datasets, covering questions from elementary to university-level. For English, we use GSM8K \citep{gsm8k}, MATH \citep{MATH}, SAT \citep{llemma}, OCW Courses \citep{minerva}, and MMLU-STEM \citep{mmlu}. Chinese datasets evaluated include MGSM-zh \citep{mgsm}, CMATH \citep{wei2023cmath}, Gaokao-MathCloze \citep{agieval}, and Gaokao-MathQA \citep{agieval}. We assess model proficiency both in generating complete text-based solutions and in solving problems via Python code, without external tools.

\paragraph{Insight on Gradient Coefficient} The methodology translates the reward signal into a gradient coefficient to drive parameter updates. In our study, we partition the reward function into `Rule'—which bases assessment on the correctness of an answer—and `Model'—which uses a trained reward model, constructed using rule-based evaluations, to assign scores to each response. Equations \ref{eq:GC-RFT} and \ref{eq:GC-GRPO} underscore an important distinction between GRPO and Online RFT: GRPO modulates the gradient coefficient adaptively according to the value produced by the reward model, thus enabling varying degrees of reinforcement or penalization depending on response quality. Conversely, Online RFT does not possess this adaptive mechanism; it provides uniform reinforcement to all correct responses and does not decrease the update magnitude for incorrect answers.

To prevent potential benchmark contamination, we adopt the filtering protocol described in \cite{deepseek-coder}. Specifically, we remove from our mathematical training set any web pages that contain questions or answers from major mathematical benchmarks in English, such as GSM8K \citep{gsm8k} and MATH \citep{MATH}, as well as from Chinese benchmarks including CMATH \citep{wei2023cmath} and AGIEval \citep{agieval}. Our filtering process operates as follows: any text passage that contains a contiguous sequence of 10 words matching any substring from these benchmarks is excluded from the corpus. For evaluation samples shorter than 10 grams but comprising at least 3 grams, we apply strict exact matching to identify and eliminate contaminated pages.

The elements comprising these approaches are outlined in Table \ref{tab:data_policy}. For a comprehensive explanation of the derivation steps, see Appendix \ref{app:analysis-rl}.

    belonging to the GLM family \citep{glm}.
\end{itemize}

Our RL experiments use DeepSeekMath-Instruct 7B as the foundation. For training, we select chain-of-thought style questions related to GSM8K and MATH from the SFT dataset, yielding a total of approximately 144K instances. To specifically assess the influence of RL on benchmarks lacking training data, we omit the remaining SFT examples from the RL phase. The reward model’s training set is built as outlined in \citep{wang2023math}. The initial reward model is trained on DeepSeekMath-Base 7B with a learning rate of 2e-5. The policy model's learning rate during GRPO is fixed at 1e-6, while the KL coefficient is set to 0.04. For every prompt, $64$ completions are sampled. The context length is capped at 1024 tokens, with a training batch size of 1024. Each exploration step leads to a single policy model update. Evaluation for DeepSeekMath-RL 7B follows the same protocol as DeepSeekMath-Instruct 7B. Tasks such as GSM8K and MATH with chain-of-thought annotations are classified as in-domain, while all remaining benchmarks are considered out-of-domain.

\begin{equation}
\footnotesize
\begin{split}
    \mathcal{J}_{GRPO}(\theta) &= \mathbb{E}{[q \sim P(Q), \{o_i\}_{i=1}^G \sim \pi_{\theta_{old}}(O|q)]}  \\
    & \frac{1}{G}\sum_{i=1}^G\frac{1}{|o_i|} \sum_{t=1}^{|o_i|} \left\{ \min \left[ \frac{\pi_\theta(o_{i,t} | q, o_{i,<t})}{\pi_{\theta_{old}}(o_{i,t} | q, o_{i,<t})} \hat{A}_{i,t}, \text{clip} \left( \frac{\pi_\theta(o_{i,t} | q, o_{i,<t})}{\pi_{\theta_{old}}(o_{i,t} | q, o_{i,<t})}, 1 - \varepsilon, 1 + \varepsilon \right)  \hat{A}_{i,t} \right] - \beta \mathbb{D}_{KL}\left[\pi_{\theta} || \pi_{ref}\right]\right\} ,
\end{split}
\label{eq:GRPO-obj}
\end{equation}

\begin{itemize}[topsep=0pt]
    \item \textbf{Closed-source models} evaluated include: (1) the GPT series, with GPT-4 \citep{gpt4} and GPT-4 Code Interpreter~\footnote{\url{https://openai.com/blog/chatgpt-plugins\#\#code-interpreter}} representing the top-tier systems, (2) Gemini Ultra and Gemini Pro \citep{gemini}, (3) Inflection-2 \citep{inflection-2}, (4) Grok-1~\footnote{\url{https://x.ai/model-card}}, as well as (5) Baichuan-3~\footnote{\url{https://www.baichuan-ai.com}}, and (6) the newest GLM-4~\footnote{\url{https://open.bigmodel.cn/dev/api\#glm-4}}

The emergence of large language models (LLM) has fundamentally transformed artificial intelligence approaches to mathematical reasoning, driving rapid progress in quantitative reasoning benchmarks \citep{MATH} as well as challenges involving geometric reasoning \citep{trinh2024solving}. Additionally, LLMs have become valuable tools for aiding humans with advanced mathematical problem-solving \citep{tao}. Yet, despite these advancements, state-of-the-art systems such as GPT-4 \citep{gpt4} and Gemini-Ultra \citep{gemini} remain inaccessible to the public, and available open-source alternatives lag considerably in their performance capabilities.

Here, we detail the procedure for tuning DeepSeekMath-Instruct 7B with mathematical instruction data, starting from the DeepSeekMath-Base model. During training, several samples are concatenated in random order until they collectively reach a total length of 4,000 tokens. The fine-tuning is conducted for 500 steps using a batch size of 256, with the learning rate fixed at 5e-5.

\textbf{The DeepSeekMath~Corpus offers exceptional quality, comprehensive multilingual mathematical coverage, and leads in dataset scale.}

\begin{algorithm}[t]
  \small
  \caption{Iterative Group Relative Policy Optimization}
  \textbf{Input} initial policy model $\pi_{\theta_{\text{init}}}$; reward models $r_\varphi$; task prompts $\mathcal{D}$; 
  hyperparameters $\varepsilon$, $\beta$, $\mu$
  \begin{algorithmic}[1]
    \State Set $\pi_\theta$ to initial policy $\pi_{\theta_{\text{init}}}$
    \For{iteration = 1, \dots, I}
       \State Define $\pi_{ref}$ as current policy $\pi_{\theta}$
      \For{step = 1, \dots, M}
      \State Draw mini-batch $\mathcal{D}_b$ from dataset $\mathcal{D}$
      \State Store previous policy $\pi_{\theta_{old}} \leftarrow \pi_{\theta}$ 
      \State For each prompt $q$ in $\mathcal{D}_b$, generate $G$ outputs $\{o_i\}_{i=1}^G \sim \pi_{\theta_{old}} (\cdot \mid q)$
      \State Use $r_{\varphi}$ to obtain rewards $\{r_i\}_{i=1}^G$ for generated outputs
      \State Compute the group relative advantage $\hat{A}_{i,t}$ at each token position $t$ for $o_i$
      \For{GRPO iteration = 1, \dots, $\mu$}
        \State Update $\pi_\theta$ by maximizing the GRPO objective (Equation \ref{eq:GC-GRPO})
      \EndFor
    \EndFor \State Continuously train $r_\varphi$ using a replay-based update scheme.
    \EndFor 
  \end{algorithmic}
  \textbf{Output} $\pi_\theta$
  \label{alg:iter-grpo}
\end{algorithm}

\subsection{Insights of Reinforcement Learning}
\subsubsection{Toward a Unified Training Framework} This section develops a coherent framework for comparing various training algorithms—namely, SFT, RFT, DPO, PPO, and GRPO—and further investigates key components within this paradigm through empirical study. The gradient with respect to the parameter $\theta$ for a given training algorithm can be formalized as:

DeepSeekMath-Base is assessed on its aptitude for mathematical problem solving through the use of few-shot chain-of-thought prompting \citep{cot}, evaluated across eight different benchmarks in both English and Chinese. These benchmarks include tests of quantitative reasoning (for example, GSM8K \citep{gsm8k}, MATH \citep{MATH}, and CMATH \citep{wei2023cmath}) as well as multiple-choice formats (such as MMLU-STEM \citep{mmlu} and Gaokao-MathQA \citep{agieval}), spanning a comprehensive array of mathematical topics ranging from basic arithmetic to advanced collegiate problems.

We begin by discussing key observations made during pre-training. Unless explicitly noted, we follow the training protocols established in Section \ref{sec:quality-policy}. Additionally, throughout this section, the term DeepSeekMath Corpus refers specifically to the 89B-token dataset collected during the second iteration of the data gathering workflow.

\item \textbf{Formal Mathematics}: DeepSeekMath-Base is assessed using the informal-to-formal theorem proving benchmark from \citep{dsp_proof} on the miniF2F dataset \citep{minif2f}, utilizing Isabelle \citep{isabelle} as the designated proof assistant. The model displays robust few-shot performance in autoformalization tasks.

Subsequent to pre-training, we perform mathematical instruction tuning for DeepSeekMath-Base, leveraging data involving chain-of-thought \citep{cot}, program-of-thought \citep{pot,pal}, and tool-integrated reasoning approaches \citep{tora}. The resultant model, DeepSeekMath-Instruct 7B, surpasses all other 7B-sized models and achieves performance levels on par with much larger 70B open-source instruction-tuned systems.

\paragraph{Insight on Data Source} We classify data sources into two types: online sampling and offline sampling. Online sampling implies that the dataset is generated dynamically by the ongoing training policy model’s exploration, whereas offline sampling utilizes data obtained from the original SFT model’s outputs. RFT and DPO utilize offline sampling, while Online RFT and GRPO rely on online sampling methods.

Moreover, we propose Group Relative Policy Optimization (GRPO), a novel reinforcement learning variant derived from Proximal Policy Optimization (PPO) \citep{schulman2017proximal}. Unlike traditional approaches that rely on a critic network, GRPO estimates the baseline using group-wise scores, dramatically reducing computational demands. Using just a select portion of English instruction tuning data, GRPO delivers significant performance gains over the already strong DeepSeekMath-Instruct, evidenced by marked improvements on both in-domain datasets (GSM8K: 82.9\% $\rightarrow$ 88.2\%, MATH: 46.8\% $\rightarrow$ 51.7\%) and out-of-domain mathematical tasks (such as CMATH: 84.6\% $\rightarrow$ 88.8\%) throughout the reinforcement learning stage. 